\section{Kiến thức nền}
\label{sec:background}

Chương này trình bày các khái niệm nền tảng cần thiết để phân tích
jailbreak attacks, jailbreak defenses và phương pháp đánh giá độ an toàn
của các mô hình ngôn ngữ lớn. Nội dung tập trung vào cách LLM sinh phản
hồi, quá trình căn chỉnh mô hình, thứ bậc chỉ dẫn, threat model, các khái
niệm an toàn liên quan và sự khác biệt giữa chatbot với agentic systems.

\subsection{Mô hình ngôn ngữ lớn}
\label{subsec:large-language-models}

Mô hình ngôn ngữ lớn (Large Language Model -- LLM) là mô hình học máy
được huấn luyện trên lượng lớn dữ liệu văn bản để học phân phối xác suất
của ngôn ngữ. Với một chuỗi token đầu vào

\begin{equation}
    x = (x_1, x_2, \ldots, x_n),
\end{equation}

mô hình tự hồi quy ước lượng xác suất của token tiếp theo dựa trên các
token đã xuất hiện trước đó:

\begin{equation}
    P_{\theta}(x_{n+1} \mid x_1, x_2, \ldots, x_n),
    \label{eq:next-token-probability}
\end{equation}

trong đó $\theta$ là tập tham số đã học của mô hình. Trong quá trình suy
luận, mô hình sinh phản hồi bằng cách dự đoán token kế tiếp, đưa token
vừa sinh trở lại ngữ cảnh, rồi tiếp tục quá trình này cho đến khi hoàn
thành phản hồi.

Trước khi được đưa vào mô hình, văn bản được chuyển thành chuỗi token
thông qua tokenizer. Nếu $s$ là chuỗi văn bản ban đầu và $\tau$ là
tokenizer, quá trình mã hóa có thể được biểu diễn bởi

\begin{equation}
    (x_1, x_2, \ldots, x_n) = \tau(s).
\end{equation}

Tokenization có vai trò quan trọng trong nghiên cứu jailbreak. Hai chuỗi
văn bản nhìn giống nhau đối với con người có thể tạo ra chuỗi token khác
nhau đối với mô hình. Ngược lại, một chuỗi token không tự nhiên hoặc khó
đọc vẫn có thể làm thay đổi mạnh hành vi của mô hình. Điều này là một cơ
sở cho các phương pháp tối ưu adversarial suffix, chẳng hạn Greedy
Coordinate Gradient (GCG) \cite{zou2023universal}.

Phần lớn LLM hiện đại được xây dựng dựa trên kiến trúc Transformer. Trong
phạm vi báo cáo này, chi tiết kiến trúc không phải trọng tâm chính. Điểm
quan trọng là LLM xử lý system prompt, chỉ dẫn của người dùng, lịch sử hội
thoại và dữ liệu bổ sung trong cùng một context window. Vì vậy, khi các
nguồn thông tin này chứa chỉ dẫn xung đột, mô hình phải dựa vào quá trình
căn chỉnh đã học để xác định hành vi phù hợp, thay vì có một cơ chế bảo
mật hình thức tuyệt đối.

\subsection{Instruction tuning và safety alignment}
\label{subsec:model-training-alignment}

Một LLM được triển khai cho người dùng thường trải qua nhiều giai đoạn
huấn luyện và căn chỉnh. Các giai đoạn cụ thể có thể khác nhau giữa từng
hệ thống, nhưng thường gồm pretraining, supervised fine-tuning,
preference alignment và safety alignment.

Trong giai đoạn \textbf{pretraining}, mô hình học dự đoán token tiếp theo
trên tập dữ liệu văn bản lớn. Giai đoạn này giúp mô hình học cấu trúc ngôn
ngữ, kiến thức tổng quát và nhiều mẫu suy luận xuất hiện trong dữ liệu.
Tuy nhiên, pretraining không trực tiếp bảo đảm mô hình sẽ làm theo chỉ dẫn
hoặc tuân thủ chính sách an toàn.

Trong giai đoạn \textbf{supervised fine-tuning} (SFT), mô hình được huấn
luyện trên các cặp dữ liệu gồm instruction và phản hồi mẫu:

\begin{equation}
    \mathcal{D}_{\mathrm{SFT}}
    =
    \left\{
        (x^{(i)}, y^{(i)})
    \right\}_{i=1}^{N},
\end{equation}

trong đó $x^{(i)}$ là chỉ dẫn và $y^{(i)}$ là phản hồi mong muốn. SFT giúp
mô hình học cách phản hồi theo định dạng hội thoại, làm theo yêu cầu và
xử lý một số tình huống phổ biến.

Trong giai đoạn \textbf{preference alignment}, mô hình được điều chỉnh dựa
trên thông tin về phản hồi nào được ưu tiên hơn. Một mẫu dữ liệu preference
có thể có dạng

\begin{equation}
    (x, y^{+}, y^{-}),
\end{equation}

trong đó $y^{+}$ là phản hồi được ưu tiên và $y^{-}$ là phản hồi kém phù
hợp hơn. Reinforcement Learning from Human Feedback (RLHF) và các phương
pháp preference optimization trực tiếp đều thuộc nhóm kỹ thuật này.

\textbf{Safety alignment} là quá trình điều chỉnh mô hình để nhận diện và
xử lý các yêu cầu có rủi ro. Mô hình được kỳ vọng từ chối yêu cầu không
được phép, cung cấp thông tin ở mức an toàn khi có thể, tránh tiết lộ chỉ
dẫn nội bộ và chuyển hướng sang nội dung phù hợp hơn. Tuy nhiên, safety
alignment không tạo ra bảo đảm tuyệt đối, vì không gian đầu vào ngôn ngữ
rất lớn trong khi dữ liệu căn chỉnh chỉ bao phủ một phần hữu hạn các tình
huống có thể xảy ra.

Chính khoảng cách giữa những mẫu đã được căn chỉnh trong huấn luyện và
những biến thể đầu vào xuất hiện tại thời điểm suy luận là một trong các
nguyên nhân khiến jailbreak vẫn tồn tại.

\subsection{Thứ bậc chỉ dẫn trong hệ thống LLM}
\label{subsec:instruction-hierarchy}

Ứng dụng LLM thường kết hợp nhiều loại thông tin trong cùng một ngữ cảnh.
Có thể phân biệt bốn nguồn chính:

\begin{enumerate}[leftmargin=1.2cm]
    \item \textbf{System instruction:}
    chỉ dẫn cấp hệ thống, quy định vai trò, chính sách và giới hạn tổng quát;

    \item \textbf{Application instruction:}
    chỉ dẫn do nhà phát triển ứng dụng thiết lập;

    \item \textbf{User instruction:}
    yêu cầu trực tiếp của người dùng;

    \item \textbf{External content:}
    nội dung từ website, tài liệu, email, cơ sở dữ liệu hoặc kết quả trả về
    từ công cụ.
\end{enumerate}

Về nguyên tắc, chỉ dẫn có độ ưu tiên cao hơn phải chi phối chỉ dẫn có độ
ưu tiên thấp hơn. Đồng thời, dữ liệu bên ngoài cần được xem là dữ liệu để
xử lý, không mặc định là chỉ dẫn điều khiển hệ thống.

\begin{figure}[H]
\centering
\begin{tikzpicture}[
    level/.style={
        draw,
        rounded corners=3pt,
        align=center,
        text width=10.5cm,
        minimum height=0.9cm,
        inner sep=5pt
    },
    >=Latex
]

\node[level] (system) {
    \textbf{Mức 1: System instruction}
    \hspace{0.4cm}
    Chính sách và giới hạn cấp hệ thống
};

\node[level, below=0.28cm of system] (developer) {
    \textbf{Mức 2: Application/Developer instruction}
    \hspace{0.4cm}
    Logic của ứng dụng
};

\node[level, below=0.28cm of developer] (user) {
    \textbf{Mức 3: User instruction}
    \hspace{0.4cm}
    Yêu cầu trực tiếp của người dùng
};

\node[level, dashed, below=0.28cm of user] (external) {
    \textbf{External content}
    \hspace{0.4cm}
    Dữ liệu không đáng tin cậy, không mặc định là instruction
};

\draw[->, thick]
    ($(system.north west)+(-0.45cm,0)$)
    --
    ($(external.south west)+(-0.45cm,0)$)
    node[
        midway,
        left=0.18cm,
        rotate=90,
        anchor=center
    ] {
        Độ ưu tiên giảm dần
    };

\end{tikzpicture}
\caption{Thứ bậc chỉ dẫn và vị trí của nội dung bên ngoài trong ứng dụng LLM.}
\label{fig:instruction-hierarchy}
\end{figure}

Trong thực tế, LLM xử lý các nguồn trên dưới dạng token trong cùng một
context window. Mô hình không có một bộ phân tích cú pháp bảo mật tuyệt đối
để phân biệt instruction với data. Vì vậy, nội dung bên ngoài có thể chứa
câu lệnh làm mô hình hiểu sai vai trò của dữ liệu. Đây là nền tảng của
prompt injection, đặc biệt là indirect prompt injection trong các hệ thống
retrieval và tool-integrated agents.

\subsection{Threat model trong nghiên cứu jailbreak}
\label{subsec:jailbreak-threat-model}

Threat model mô tả mục tiêu, quyền truy cập, kiến thức và giới hạn của kẻ
tấn công. Việc xác định threat model là cần thiết vì hai phương pháp
jailbreak không thể được so sánh công bằng nếu chúng dựa trên các giả định
khác nhau.

Bảng~\ref{tab:jailbreak-threat-model-dimensions} tổng hợp các chiều thường
được sử dụng trong nghiên cứu jailbreak.

\begin{table}[H]
\centering
\caption{Các chiều chính của threat model trong nghiên cứu jailbreak.}
\label{tab:jailbreak-threat-model-dimensions}
\begin{tabularx}{\textwidth}{
    >{\bfseries}p{0.23\textwidth}
    p{0.29\textwidth}
    X
}
\toprule
Chiều đánh giá
& Các trường hợp phổ biến
& Ý nghĩa \\
\midrule

Quyền truy cập
& White-box, gray-box, black-box
& Attacker có thể truy cập tham số, gradient, logit hoặc chỉ giao tiếp
qua đầu vào--đầu ra. \\

Số lượt tương tác
& Single-turn, multi-turn
& Attack được thực hiện trong một truy vấn hoặc qua nhiều lượt hội thoại. \\

Khả năng thích nghi
& Static, adaptive
& Attacker có thay đổi chiến lược dựa trên phản hồi hoặc defense hay không. \\

Ngân sách truy vấn
& Giới hạn hoặc không giới hạn
& Phản ánh chi phí và tính khả thi của attack, đặc biệt trong black-box
setting. \\

Kiến thức về hệ thống
& Không biết, biết một phần, biết đầy đủ
& Attacker có biết model family, system prompt, policy hoặc defense đang
được sử dụng hay không. \\

Mức tự động hóa
& Manual, template-based, optimization-based
& Xác định mức độ phụ thuộc vào con người và thuật toán tìm kiếm. \\

Mục tiêu
& Unsafe response, policy bypass, tool misuse
& Hành vi mà attacker muốn hệ thống thực hiện. \\

Đối tượng
& Chatbot, multimodal model, agent, robot
& Quyết định phạm vi hậu quả và loại evaluator cần sử dụng. \\

Transferability
& Model-specific, transferable
& Attack chỉ hiệu quả trên một mô hình hay có thể chuyển sang mô hình khác. \\

\bottomrule
\end{tabularx}
\end{table}

Trong các chương sau, các chiều này được dùng để phân loại những nhóm tấn
công khác nhau. Ví dụ, white-box attacks có thể tận dụng thông tin nội bộ
của mô hình để tối ưu đầu vào, trong khi black-box attacks chỉ dựa trên
việc gửi prompt và quan sát phản hồi. Tương tự, adaptive attacks thường
khó phòng thủ hơn static attacks vì attacker có thể điều chỉnh chiến lược
dựa trên cơ chế defense.

\subsection{Phân biệt các khái niệm an toàn liên quan}
\label{subsec:related-security-concepts}

Các thuật ngữ jailbreaking, prompt injection, adversarial example, data
poisoning và backdoor thường được sử dụng gần nhau nhưng mô tả các threat
model khác nhau. Việc phân biệt các khái niệm này giúp xác định đúng phạm
vi của báo cáo.

\begin{table}[H]
\centering
\caption{Phân biệt jailbreaking với các khái niệm an toàn liên quan.}
\label{tab:security-concept-comparison}
\begin{tabularx}{\textwidth}{
    >{\bfseries}p{0.20\textwidth}
    p{0.25\textwidth}
    p{0.25\textwidth}
    X
}
\toprule
Khái niệm
& Thời điểm tác động
& Mục tiêu chính
& Ví dụ bề mặt tác động \\
\midrule

Jailbreaking
& Thời điểm suy luận
& Vượt qua ràng buộc an toàn của mô hình
& User prompt, suffix hoặc hội thoại nhiều lượt \\

Prompt injection
& Thời điểm suy luận
& Thay đổi luồng điều khiển hoặc thứ tự ưu tiên chỉ dẫn
& User prompt hoặc dữ liệu được đưa vào context \\

Indirect prompt injection
& Thời điểm suy luận
& Làm agent chịu ảnh hưởng bởi chỉ dẫn nằm trong nguồn dữ liệu ngoài
& Website, email, tài liệu hoặc tool output \\

Adversarial example
& Thời điểm suy luận
& Làm mô hình dự đoán hoặc hành xử sai dưới nhiễu có chủ đích
& Token, hình ảnh, âm thanh hoặc embedding \\

Data poisoning
& Giai đoạn huấn luyện
& Làm thay đổi hành vi mô hình thông qua dữ liệu bị nhiễm
& Pretraining corpus hoặc fine-tuning dataset \\

Backdoor
& Huấn luyện hoặc chỉnh sửa mô hình
& Kích hoạt hành vi ẩn khi xuất hiện trigger
& Trigger token, mẫu văn bản hoặc đặc trưng đầu vào \\

\bottomrule
\end{tabularx}
\end{table}

Jailbreaking chủ yếu nhằm làm mô hình vi phạm chính sách an toàn. Prompt
injection tập trung vào việc khiến ứng dụng bỏ qua hoặc thay đổi chỉ dẫn
ban đầu. Trong chatbot đơn giản, hai khái niệm này có thể giao nhau. Trong
agentic system, prompt injection có thể không yêu cầu mô hình tạo nội dung
không an toàn, mà chỉ cần làm agent thực hiện một công cụ, truy xuất dữ
liệu hoặc thay đổi kế hoạch ngoài ý định của người dùng.

\subsection{Safety--utility trade-off}
\label{subsec:safety-utility-tradeoff}

Một hệ thống an toàn không chỉ cần giảm phản hồi không phù hợp mà còn phải
duy trì khả năng hỗ trợ các yêu cầu hợp lệ. Nếu chỉ tối ưu để mô hình từ
chối càng nhiều càng tốt, hệ thống có thể đạt tỷ lệ unsafe compliance thấp
nhưng gần như không còn giá trị sử dụng.

Có thể phân biệt hai lỗi quan trọng:

\begin{enumerate}[leftmargin=1.2cm]
    \item \textbf{Unsafe compliance:}
    mô hình chấp nhận hoặc thực hiện yêu cầu không an toàn. Đây là kết quả
    mà jailbreak attack thường hướng tới.

    \item \textbf{Over-refusal:}
    mô hình từ chối một yêu cầu hợp lệ. Đây là tác dụng phụ thường gặp khi
    defense quá thận trọng.
\end{enumerate}

Vì vậy, đánh giá một defense không nên chỉ dựa trên Attack Success Rate.
Một phương pháp phòng thủ đầy đủ cần được kiểm tra trên cả tập harmful
prompts để đánh giá safety và tập benign prompts để đánh giá utility. Ngoài
ra, cần xem xét chi phí suy luận, độ trễ, số lần gọi mô hình và khả năng
chống lại adaptive attacker.

Sự đánh đổi giữa safety và utility là một chủ đề xuyên suốt của báo cáo:
defense quá yếu làm tăng rủi ro jailbreak, trong khi defense quá nghiêm
ngặt làm tăng over-refusal và làm giảm trải nghiệm người dùng.

\subsection{Từ chatbot đến agentic systems}
\label{subsec:from-chatbots-to-agents}

Chatbot dựa trên LLM thường nhận yêu cầu và trả về văn bản. Agentic system
mở rộng mô hình này bằng cách bổ sung các thành phần như bộ nhớ, bộ lập kế
hoạch, retrieval từ nguồn dữ liệu bên ngoài, khả năng gọi công cụ, quan sát
trạng thái môi trường và vòng lặp hành động--phản hồi.

Có thể biểu diễn một trajectory của agent dưới dạng

\begin{equation}
    \tau
    =
    \left(
        o_0,
        a_0,
        o_1,
        a_1,
        \ldots,
        o_T
    \right),
    \label{eq:agent-trajectory}
\end{equation}

trong đó $o_t$ là observation tại thời điểm $t$ và $a_t$ là hành động được
agent lựa chọn.

Đối với chatbot, evaluator thường tập trung vào phản hồi cuối. Đối với
agentic system, đánh giá cần xem xét cả trajectory, bao gồm kế hoạch, lời
gọi công cụ, trạng thái trung gian và hành động cuối cùng. Một agent có
thể tạo ra phản hồi văn bản an toàn nhưng vẫn thực hiện lời gọi công cụ
không phù hợp ở bước trung gian. Ngược lại, một kế hoạch không an toàn có
thể bị action controller chặn trước khi hành động được thực thi.

\begin{table}[H]
\centering
\caption{So sánh bề mặt an toàn của chatbot và agentic system.}
\label{tab:chatbot-agent-security-comparison}
\begin{tabularx}{\textwidth}{
    >{\bfseries}p{0.22\textwidth}
    p{0.32\textwidth}
    X
}
\toprule
Khía cạnh
& Chatbot
& Agentic system \\
\midrule

Đầu vào
& User prompt và lịch sử hội thoại
& User prompt, memory, retrieval, tool output và environment \\

Đầu ra
& Chủ yếu là văn bản
& Văn bản, kế hoạch, lời gọi công cụ và hành động \\

Đơn vị đánh giá
& Một phản hồi hoặc một hội thoại
& Toàn bộ trajectory của agent \\

Bề mặt tấn công
& Jailbreak và direct prompt injection
& Jailbreak, indirect injection, memory poisoning và tool manipulation \\

Hậu quả
& Nội dung phản hồi không phù hợp
& Có thể ảnh hưởng dữ liệu, dịch vụ hoặc môi trường vật lý \\

Phòng thủ
& Input/output guard và alignment
& Guard, sandbox, permission control, monitoring và human confirmation \\

\bottomrule
\end{tabularx}
\end{table}

Sự mở rộng từ chatbot sang agentic systems là lý do báo cáo không chỉ phân
tích jailbreak trên LLM sinh văn bản, mà còn xem xét các hệ thống có khả
năng sử dụng công cụ và thực hiện hành động. Với các hệ thống này, an toàn
không thể chỉ phụ thuộc vào hành vi từ chối của mô hình, mà cần được thiết
kế ở cấp kiến trúc.
